\chapter{MLIR: Modern Compiler Infrastructure}
\label{chap:ch14}
\typeout{START_CHAPTER "ch14" \theabspage}

The slide deck says ``one IR, every target.'' The production dashboard says something else: the same PyTorch-exported decoder graph, lowered through MLIR, can differ by an order of magnitude in p99 latency across GPU, CPU, and NPU fleets \citep{patel2025llminference,chen2020compilerbasedhardw}. That gap is rarely a missing dialect---it is a missing \emph{contract}. Chapter~3 mapped residency and legality per hardware class; this chapter shows how MLIR must encode those predicates in IR attributes, pass gates, and emit paths before any canonicalizer runs \citep{amd2024xdna,yirage2026}. Every section below follows Chapter~1's discipline: name the decode failure mode, trace it to a hardware tier, show the pass or emit fix, and close with bytes/token or launches/token---not peak compile-time FLOPs \citep{taxbreak2026,memoryfloor2026}.

TaxBreak's decode tables are the mental anchor. Dense Llama-3.2-1B at BS${=}4$/SL${=}2048$ averages 847.5 kernel launches per output token on H100; OLMoE-1B/7B averages 9,305 \citep{taxbreak2026}. MLIR does not erase those numbers by itself---it \emph{enables} passes and backends to collapse launch DAGs when residency metadata survives bufferization \citep{nvidia2024cudagraphs,dao2022flashattention}. CUDA Graphs on Qwen-2.5-7B (BS${=}1$, ctx${=}2048$) removes ${\approx}3.05$\,ms of host overhead per step on H100 ($1.259\times$ median speedup) yet does nothing for illegal HBM traffic introduced by bad bufferization \citep{memoryfloor2026}. The five compiler sections below---hardware matrix, passes, codegen, tuning, case study---mirror the template used for XLA, TVM, and Triton so Part~VI comparisons stay metric-aligned \citep{chen2020compilerbasedhardw}.

\section{MLIR HW matrix}
\label{sec:ch14_mlir_hw_matrix}

\paragraph{Why decode breaks the ``portable IR'' story.} Prefill hides target sins: large $L$ amortizes launch overhead and keeps tensor cores fed \citep{dao2022flashattention,patel2025llminference}. Decode removes that amortization---batch-1, one token, full weight + KV read every step \citep{memoryfloor2026,dlbook}. MLIR modules that legalize during prefill tuning often fail decode legality checks because dynamic dimensions tied to sequence length propagate into bufferization \citep{kwon2023vllm,amd2024xdna}. The hardware matrix must therefore include a decode row, not only a ``training/inference'' row \citep{yirage2026,taxbreak2026}.

\paragraph{Worked numbers.} TaxBreak's 847.5 launches/token (dense Llama-3.2-1B) and 9,305 launches/token (OLMoE) bound what MLIR fusion must collapse before MMA micro-optimizations matter \citep{taxbreak2026}. Memory-Floor's $1.259\times$ Graphs speedup on Qwen-2.5-7B shows host-side wins when launch DAGs stabilize---but Graphs cannot fix alloc-heavy IR \citep{memoryfloor2026,nvidia2024cudagraphs}. Matrix gates should reference both metrics \citep{yirage2026,chen2020compilerbasedhardw}.
% AUTO_VISUAL:fig_mlir_hw_matrix_pipeline

\begin{figure}[t]
\centering
\begin{tikzpicture}[vis base, scale=0.88, transform shape]
 \node[vis pipeline node] (s0) at (-5.03,0) {\footnotesize Input\\MLIR HW Matrix};
 \node[vis arch node] (s1) at (-1.68,0) {\footnotesize On-chip\\Tiling / residency};
 \node[vis pipeline node] (s2) at (1.67,0) {\footnotesize Fused compute\\MegaKernel stage};
 \node[vis arch node] (s3) at (5.03,0) {\footnotesize Output\\Next layer / KV};
 \draw[vis arrow] (s0.east) -- (s1.west);
 \draw[vis arrow] (s1.east) -- (s2.west);
 \draw[vis arrow] (s2.east) -- (s3.west);
\end{tikzpicture}
\caption{MLIR target matrix: the same Linalg matmul must lower through different memory tiers per target.}
\label{fig:mlir_hw_matrix_pipeline}
\end{figure}

The recurring mistake is publishing a target checklist---GPU, CPU, NPU all checked---without stating which \emph{hardware} features each backend actually models \citep{nvidia2022hopper,amd2024xdna}. A row that says ``NPU: supported'' is meaningless if the lowering path assumes dynamic tensor shapes and host-side allocators that void static buffer tables on XDNA \citep{amd2024aie}. For decode, the matrix must speak residency: shared-memory caps, DMA descriptor limits, cache blocking rules, and whether tensor-core MMA shapes are mandatory or optional \citep{semianalysis2024tma}.

\begin{table}[t]
\centering
\caption{MLIR backend capability matrix for LLM decode (conceptual).}
\label{tab:mlir_hw_matrix}
\begin{tabular}{lccc}
\toprule
\textbf{Dimension} & \textbf{GPU (NVVM/ROCDL)} & \textbf{CPU (LLVM)} & \textbf{NPU (XDNA/AIE)} \\
\midrule
Dynamic batch/seq & Flexible & Flexible & Restricted \\
On-chip scratch & Shared mem / L1 & L2/L3 tiles & Tile SRAM \\
Fusion default & Partial (Linalg) & Cache-blocked & Full static \\
Launch model & SIMT + graphs & OpenMP/SIMD & Spatial DMA \\
Decode pain & Launches/token & LLC misses & Static legality \\
\bottomrule
\end{tabular}
\end{table}

YiRage treats Table~\ref{tab:mlir_hw_matrix} as typed predicates in its MLIR importers: a pass that introduces a dynamic dimension on an NPU row is rejected before codegen, not at runtime \citep{yirage2026}. GPU rows must declare TMA/async-copy legality; CPU rows must attach NUMA hints when KV shards cross sockets \citep{kwon2023vllm,dlbook}. The matrix is a \emph{gate}, not documentation---identical to Chapter~3's constraint matrix blocking illegal MegaKernel fusion early \citep{dao2022flashattention}.

Dialect designers often ship one Linalg op set for matmul, softmax, and layer norm, then expect backends to infer residency \citep{chen2020compilerbasedhardw}. Decode exposes the gap: GPU backends need explicit memory-space attributes on tensors headed for shared memory; CPU backends need layout markers for cache-line alignment; NPU backends need lifetime intervals tied to DMA events \citep{nvidia2022hopper,amd2024xdna}. Without those attributes, the matrix devolves into marketing---every target checked, none explained \citep{yirage2026}.

JIT compilation amplifies the risk. MLIR JIT can specialize tile sizes at runtime for dynamic sequence lengths---valuable for prefill---but destructive for decode when specialization reintroduces host-side planning per token \citep{taxbreak2026,nvidia2024cudagraphs}. YiRage restricts JIT to shape buckets that preserve static launch DAGs on GPU and fully static buffer tables on NPU; everything else triggers recompilation with explicit user-visible cost \citep{yirage2026,patel2025llminference}.

\paragraph{Target interfaces vs.\ checkmarks.} MLIR conversion interfaces are where the matrix becomes executable \citep{chen2020compilerbasedhardw}. Backends advertise supported ops, memory spaces, and layout constraints; importers must attach attributes \emph{before} canonicalization erases them \citep{yirage2026}. Decode teams should reject merges when a matrix row flips from supported to unsupported without a residency annotation change---silent downgrades are how bytes/token regressions enter production \citep{memoryfloor2026}.

\paragraph{ROCm and CUDA divergence.} NVVM and ROCDL lowerings diverge on warp shuffle availability, LDS sizing, and MFMA vs WMMA shapes \citep{nvidia2022hopper,liu2024deepseekv3}. An MLIR module tuned for Hopper TMA may compile on MI300 yet emit extra buffer reloads because async-copy attributes do not map one-to-one \citep{semianalysis2024tma}. YiRage keeps separate predicate tables per vendor because decode SLOs are vendor-local even when Linalg looks identical \citep{yirage2026,taxbreak2026}.

\begin{example}
Hardware matrix gate: an importer adds \texttt{memref.alloc} on the decode hot path for a softmax statistic. GPU row passes (HBM spill is slow but legal); NPU row fails static legality; CPU row passes but raises bytes/token in Memory-Floor sweeps \citep{memoryfloor2026,amd2024xdna}. The matrix should have blocked the NPU alloc and flagged GPU/CPU regressions in CI \citep{yirage2026}.
\end{example}

\section{MLIR dialect stack for decode export}
\label{sec:ch14_dialect_stack}

\paragraph{Export pitfalls.} Torch-MLIR and similar exporters often emit verbose \texttt{tensor} graphs with redundant casts \citep{chen2020compilerbasedhardw,yirage2026}. Each cast is a future buffer unless erased before bufferization \citep{memoryfloor2026}. Decode exporters should run target-tagged cleanup passes \emph{before} generic canonicalization \citep{amd2024xdna,patel2025llminference}. Skipping cleanup yields IR that looks portable yet allocates like eager PyTorch \citep{taxbreak2026,kwon2023vllm}.

Exporting a decoder block from PyTorch lands in \texttt{func.func} wrappers around \texttt{linalg} ops on \texttt{tensor} values, then converts to \texttt{memref} through bufferization \citep{chen2020compilerbasedhardw,yirage2026}. Decode-sensitive details live in the transitions: \texttt{tensor} world allows implicit temporaries; \texttt{memref} world exposes every alloc \citep{memoryfloor2026}. Teams that benchmark only pre-export \texttt{linalg} miss the alloc explosion bufferization introduces \citep{kwon2023vllm,taxbreak2026}.

\texttt{scf.for} and \texttt{scf.if} encode loop structure for sequence dimensions \citep{dao2023flashdecoding}. For decode, outer loops over tokens should remain on the host while inner tiles map to device parallelism---but generic loop-invariant-code-motion can hoist device launches into host loops if passes lack target annotations \citep{nvidia2024cudagraphs,yirage2026}. YiRage marks host vs device regions explicitly before any LICM runs \citep{yirage2026,patel2025llminference}.

The \texttt{gpu} dialect bridges to launches; \texttt{nvvm} or \texttt{rocdl} attach machine intrinsics \citep{nvidia2022hopper}. On CPU, \texttt{vector} and \texttt{llvm} carry SIMD width choices \citep{dlbook}. On NPU, vendor dialects replace \texttt{gpu.launch} with spatial pipelines \citep{amd2024aie,amd2024xdna}. A decode compiler must never assume a single lowering path from \texttt{linalg.matmul}---the stack is a tree, not a line \citep{chen2020compilerbasedhardw}.

\begin{example}
PyTorch export produces \texttt{linalg.softmax} on a decode slice. Without residency tags, bufferization allocates a full-sequence temp; with tags tied to KV views, the same op updates in-place carriers compatible with Flash-Decoding-style parallelism \citep{dao2022flashattention,dao2023flashdecoding,yirage2026}.
\end{example}

Dialect legalization should run per target \emph{before} shared canonicalization \citep{yirage2026}. Shared canonicalizers fold ops that backends rely on to prove static shapes \citep{amd2024xdna}. Canonicalize too early and you lose the hints backends need---the LLVM lesson repeats on every MLIR upgrade \citep{chen2020compilerbasedhardw,memoryfloor2026}.

\section{MLIR arch passes}
\label{sec:ch14_mlir_arch_passes}

\paragraph{Affine maps and paged KV.} PagedAttention requires MLIR to express gather/scatter on KV blocks with provable in-bounds access \citep{kwon2023vllm,dao2022flashattention}. Lowering paged views late hides page boundaries from fusion passes---exactly how illegal cross-page fusion enters production IR \citep{yirage2026,memoryfloor2026}. Run paging-aware bufferization before \texttt{linalg-fuse} on all targets \citep{taxbreak2026,amd2024aie}.

\paragraph{MoE pass boundaries.} Expert routing introduces control flow that tempts pass writers to split launches per expert \citep{taxbreak2026,liu2024deepseekv3}. Decode MoE needs guarded fusion: batch experts to cap launches while respecting SMEM \citep{yirage2026,patel2025llminference}. Pass pipelines should expose expert-count limits as first-class attributes \citep{chen2020compilerbasedhardw,nvidia2022hopper}.
% AUTO_VISUAL:fig_mlir_arch_passes_architecture

\begin{figure}[t]
\centering
\begin{tikzpicture}[vis base, scale=0.88, transform shape]
 \node[vis arch node] (s0) at (-5.03,0) {\footnotesize Host\\};
 \node[vis arch node] (s1) at (-1.68,0) {\footnotesize On-chip buffer\\};
 \node[vis arch node] (s2) at (1.67,0) {\footnotesize Compute array\\};
 \node[vis arch node] (s3) at (5.03,0) {\footnotesize Memory tier\\};
 \draw[vis arrow] (s0.east) -- (s1.west);
 \draw[vis arrow] (s1.east) -- (s2.west);
 \draw[vis arrow] (s2.east) -- (s3.west);
\end{tikzpicture}
\caption{Pass pipeline: dialect lowering must preserve producer--consumer residency across tiers.}
\label{fig:mlir_arch_passes_architecture}
\end{figure}

MLIR's composable pass pipelines---bufferization, Linalg tiling, vectorization, dialect conversion---feel like power until pass order silently changes memory traffic \citep{chen2020compilerbasedhardw}. A canonical failure: aggressive \texttt{linalg-fuse} before bufferization on GPU lowers decode attention into a single kernel on paper, yet the lowered IR still materializes softmax statistics to HBM because the dialect contract did not keep partial sums in shared memory \citep{dao2022flashattention,dao2023flashdecoding}. The fix is not fewer passes; it is tagging each pass with the hardware dimension it optimizes (SRAM footprint, launch count, bandwidth) and the targets where it may run \citep{yirage2026}.

Cross-hardware lowering diverges at the same Linalg op. GPU lowering targets warp tiles aligned to MMA instructions; CPU lowering emits AVX/AMX micro-kernels with cache blocking; NPU lowering emits spatial schedules with fixed buffer lifetimes \citep{amd2024xdna,nvidia2022hopper}. A vectorization pass tuned for CPU can break NPU legality by introducing strided layouts incompatible with DMA alignment---regressions that appear only in bytes/token dashboards, not unit tests \citep{memoryfloor2026}. Industrial pipelines fork pass registries per target early, sharing only algebraic simplification \citep{liu2024deepseekv3}.

\begin{example}
Single-layer decode attention in Linalg: on Hopper, \texttt{lower-vector-gpu} plus shared-memory promotion can fuse QK$^\top$ and softmax carriers when tiles fit SMEM; on XDNA the same graph must be fully bufferized to static tile sizes before AIE codegen, or the backend falls back to DDR-limited scalar loops \citep{amd2024aie,yirage2026,semianalysis2024tma}.
\end{example}

Bufferization deserves decode-specific scrutiny. Generic one-shot bufferization assumes heap fallback \citep{chen2020compilerbasedhardw}. Decode attention cannot afford that: KV tensors must alias paged pools (GPU), huge-page arenas (CPU), or fixed SRAM banks (NPU) \citep{kwon2023vllm,amd2024aie}. Pass pipelines should declare allocation tiers in IR so lowering never inserts silent \texttt{memref.alloc} on the hot path \citep{yirage2026,memoryfloor2026}.

Dialect conversion ordering is another launch multiplier. Converting \texttt{linalg} $\rightarrow$ \texttt{vector} $\rightarrow$ \texttt{gpu} in the wrong order explodes a fused region into dozens of micro-kernels---the fragmentation Chapter~1 measured in eager PyTorch \citep{taxbreak2026}. Production stacks freeze a golden pass list per target and diff IR dumps when upgrading LLVM or MLIR \citep{liu2024deepseekv3,chen2020compilerbasedhardw}.

\paragraph{One-shot vs.\ iterative bufferization.} One-shot bufferization is fast to wire but assumes heap fallback; iterative bufferization can reuse buffers across fused regions yet increases compile time---acceptable for offline NPU builds, painful for GPU JIT \citep{amd2024aie,yirage2026}. Pick strategy per target; mixing strategies within one registry causes ``works in unit test, fails at BS${=}1$'' bugs \citep{memoryfloor2026,kwon2023vllm}.

\paragraph{Pass instrumentation.} Attach counters to passes: bytes introduced, allocs inserted, launch boundaries created \citep{taxbreak2026}. A pass that reduces op count but increases alloc count is suspect for decode \citep{memoryfloor2026,yirage2026}.

\section{MLIR codegen}

Teams celebrate the first successful NVVM lowering, then watch p99 decode flatline because the \emph{platform} wrapper still constructs tensors on the host every token \citep{taxbreak2026,nvidia2024cudagraphs}. Codegen is not a backend detail---it is the last place residency contracts survive or die before the driver sees them \citep{yirage2026,chen2020compilerbasedhardw}.
\label{sec:ch14_mlir_codegen}

\paragraph{Graph capture contract.} CUDA Graphs require identical launch addresses and dependency topology across replay \citep{nvidia2024cudagraphs,taxbreak2026}. MLIR emit that injects pointer-chasing host setup between device kernels breaks capture silently---graphs record yet decode latency unchanged \citep{memoryfloor2026,yirage2026}. Encode capture legality as a codegen predicate alongside occupancy \citep{patel2025llminference,chen2020compilerbasedhardw}.

\paragraph{CPU and NPU emit deltas.} CPU emit must pin threads and align KV for cache lines; NPU emit must assign DMA lifetimes to match tile SRAM \citep{dlbook,amd2024aie,amd2024xdna}. Identical high-level MLIR is fine; identical emit is not \citep{yirage2026,liu2024deepseekv3}.

Production incidents blame ``the NVVM backend'' when the real bug was emitting LLVM IR that assumes a mutable stack and per-op \texttt{malloc} in the host wrapper---poison for CUDA Graphs and NPU static schedules alike \citep{nvidia2024cudagraphs,taxbreak2026}. Codegen is where MLIR's abstraction meets launch reality: GPU emit must respect stream ordering, cluster scope, and TMA descriptor lifetimes; CPU emit must preserve NUMA-local buffers; NPU emit must serialize DMA chains with compile-time-known latencies \citep{nvidia2022hopper,amd2024xdna}.

Backend selection is not neutral. NVVM/ROCDL paths excel when decode graphs contain dynamic control flow that still maps to SIMT; bare-metal AIE paths excel when the entire decoder step is a static dataflow graph \citep{amd2024aie}. YiRage's MLIR backends share high-level Linalg but diverge at emit: CUDA/HIP paths generate MegaKernel-style fused launches; XDNA paths emit buffer descriptors and ping-pong DMA schedules \citep{yirage2026,semianalysis2024tma}. One emit artifact for all targets guarantees either GPU launch storms or NPU compile failures \citep{patel2025llminference}.

Emit quality shows up in host wrappers. MLIR-to-CUDA that allocates temporaries each launch breaks graph capture; decode pays the TaxBreak orchestration tax \citep{taxbreak2026,nvidia2024cudagraphs}. Good codegen hoists allocations to capture time or eliminates them via shared-memory arenas \citep{memoryfloor2026}. On CPU, emit must avoid false sharing in OpenMP regions around thin decode GEMMs \citep{dlbook}.

\paragraph{Host wrapper ABI.} Codegen finishes only when the host wrapper is graph-capturable. \texttt{cudaMalloc} per launch breaks capture; stack pointers into kernels break on some driver versions \citep{nvidia2024cudagraphs,memoryfloor2026}. Hoist workspace to session init; pass arena offsets as kernel arguments \citep{yirage2026}.

\paragraph{LLVM pipeline forks.} GPU paths often stop at NVVM/ROCDL; CPU paths continue through LLVM vectorizers; NPU paths may bypass LLVM entirely \citep{chen2020compilerbasedhardw,amd2024aie}. Sharing one LLVM optimization stage invites loop unrolling that expands register pressure on GPU decode tiles \citep{nvidia2022hopper,liu2024deepseekv3}.

\paragraph{Debuggability vs.\ fusion.} Aggressive fusion produces opaque dumps. Keep named MLIR anchors (\texttt{decode\_attn\_tile\_7}) through emit so perf teams map regressions to passes \citep{yirage2026,taxbreak2026}.

\section{MLIR tuning}

Profile logs from prefill runs are the classic trap: they recommend tile sizes that overflow shared memory during batch-1 decode and force HBM spill \citep{memoryfloor2026,semianalysis2024tma}. Autotuning MLIR without a decode-shaped \emph{profile} trace repeats the same mistake at compiler scale \citep{taxbreak2026,chen2020compilerbasedhardw}.
\label{sec:ch14_mlir_tuning}

\paragraph{Search spaces.} Constrain tile search by SMEM capacity, warp count, and max fusion depth derived from Chapter~3 models \citep{nvidia2022hopper,amd2024xdna,yirage2026}. Unconstrained search finds prefill winners that spill on decode \citep{memoryfloor2026,semianalysis2024tma}. Log Pareto points with launches/token on the decode trace, not FLOPs on square GEMMs \citep{taxbreak2026,patel2025llminference}.

\paragraph{Fleet replay.} Version autotune artifacts with MLIR commit, LLVM commit, and model hash \citep{chen2020compilerbasedhardw,liu2024deepseekv3}. Replay canned decode traces in CI before fleet rollout---same discipline as CUDA Graphs templates \citep{nvidia2024cudagraphs,memoryfloor2026,taxbreak2026}.

Autotuning MLIR pipelines---tile sizes, fusion depths, vector widths---looks like free performance until you ship prefill-shaped winners to batch-1 decode fleets \citep{memoryfloor2026,chen2020compilerbasedhardw}. Profile-guided passes that maximize FLOPs during autotune often increase shared-memory footprint past decode legality, forcing silent spill to HBM \citep{taxbreak2026,semianalysis2024tma}. Always co-tune with launches/token: one extra kernel boundary per layer can erase tens of microseconds per token on MoE models \citep{liu2024deepseekv3,taxbreak2026}.

The dominant autotune \emph{pitfall} on decode: optimizing square GEMMs while ignoring launch count. GPU autotune must cap warps and SMEM; CPU autotune must model LLC capacity and thread pinning; NPU autotune must not explore dynamic shapes that invalidate DMA tables \citep{amd2024xdna,dlbook}. Cross-hardware profile imports are dangerous---GPU configs thrash CPU caches; NPU may not compile at all \citep{kwon2023vllm,yirage2026}. YiRage encodes per-target search spaces so tuners cannot propose illegal configs \citep{yirage2026,patel2025llminference}.

\paragraph{Cost models for decode.} Combine estimated bytes moved, predicted launches, and SMEM footprint against hardware caps \citep{memoryfloor2026,taxbreak2026}. MLIR's transform dialect can encode constraints before search starts \citep{chen2020compilerbasedhardw}. Without them, RL tuners win on square GEMMs yet spill during fused decoder blocks \citep{yirage2026}.

\paragraph{Regression gates.} Tie autotune merges to decode traces. Require non-regression on launches/token and bytes/token at BS${=}1$, ctx${=}2048$ before accepting new tile configs \citep{taxbreak2026,memoryfloor2026,nvidia2024cudagraphs}.

\section{MLIR case study}
\label{sec:ch14_mlir_case_study}

\paragraph{Triplet harness.} YiRage-style triplet runs share tokenizer, context bucket, and batch policy \citep{yirage2026,patel2025llminference}. Report bytes/token, launches/token, compile rejections, and p99 alongside means \citep{memoryfloor2026,taxbreak2026}. A compile rejection on NPU is a first-class outcome---not a footnote \citep{amd2024xdna,amd2024aie}.

\paragraph{Silent spill class E.} Category E failures compile cleanly while bytes/token jump because SMEM promotion failed \citep{memoryfloor2026,dao2022flashattention}. Detect by diffing allocs between pass stages; fix by reordering bufferization relative to fusion \citep{yirage2026,chen2020compilerbasedhardw,taxbreak2026}.

Consider a minimal decode cell: Llama-class decoder layer, batch-1, context 2048, exported from PyTorch to MLIR via Linalg \citep{patel2025llminference}. The question is not ``does it compile?'' but ``does bytes/token match the MegaKernel baseline from Part~IV?'' On H100, an unfused MLIR pipeline may compile cleanly yet issue hundreds of launches per token---the pathology TaxBreak measures in framework stacks \citep{taxbreak2026}. Fused Linalg plus GPU codegen and CUDA Graphs capture can recover host overhead when emit preserves a static launch DAG \citep{nvidia2024cudagraphs,memoryfloor2026}.

The identical module lowered to x86 LLVM with OpenMP may achieve acceptable throughput on small models but regress on large weights due to LLC pressure \citep{dlbook}. Lowered to XDNA, the module either becomes a static MegaKernel with competitive watt/token or fails legality checks---rarely a middle ground \citep{amd2024xdna,amd2024aie}. YiRage's benchmark harness runs this triplet on every merge: same IR, three backends, report bytes/token, launches/token, and compile-time legality violations \citep{yirage2026}. That triplet is empirical proof of the book's thesis: hardware dictates compiler outcomes, not the reverse \citep{chen2020compilerbasedhardw}.

\begin{example}
Cross-hardware MLIR decode cell (Llama-3.2-1B layer, BS${=}1$, ctx${=}2048$): GPU fused+graphs targets competitive ms/token only when launches collapse toward TaxBreak dense baselines; CPU LLVM trades latency for deployment flexibility; XDNA wins only when fusion is total and DMA-static \citep{taxbreak2026,memoryfloor2026,yirage2026}.
\end{example}

\paragraph{Paged KV and MLIR types.} PagedAttention indirection must appear as first-class \texttt{memref} views with lifetime intervals---not opaque pointers lowered late \citep{kwon2023vllm,yirage2026}. Late lowering hides page boundaries and causes illegal fusion across pages \citep{dao2022flashattention,memoryfloor2026}.

\paragraph{MoE-specific notes.} MoE layers explode dispatch when each expert lowers to separate launches \citep{taxbreak2026,liu2024deepseekv3}. MLIR pipelines need expert-batch fusion guarded by SMEM caps \citep{yirage2026}. Case studies that report dense Llama only understate orchestration challenge---pair dense and MoE exports \citep{patel2025llminference,taxbreak2026}.

\begin{table}[t]
\centering
\caption{Case-study outcome taxonomy for cross-hardware MLIR decode (illustrative).}
\label{tab:mlir_case_taxonomy}
\begin{tabular}{lll}
\toprule
\textbf{Class} & \textbf{Symptom} & \textbf{Typical fix} \\
\midrule
B & High launches/token & Fuse + graph capture \\
C & High bytes/token & SMEM/DMA tiling \\
D & NPU compile fail & Static shape specialization \\
E & Flat latency, rising bytes & Fix bufferization order \\
\bottomrule
\end{tabular}
\end{table}


\paragraph{Industrial matrix CI.} Mature MLIR importers check machine-readable matrix files into CI \citep{yirage2026,chen2020compilerbasedhardw}. Each row lists supported ops, max tile sizes, and forbidden dynamic dimensions \citep{amd2024xdna,kwon2023vllm}. When PyTorch export adds rotary embedding or a new norm variant, the matrix diff tells reviewers which targets lose legality \emph{before} any benchmark runs \citep{patel2025llminference,liu2024deepseekv3}. This is how decode teams avoid green GPU CI paired with red NPU field releases \citep{memoryfloor2026,taxbreak2026}.

\paragraph{Pass-order war stories.} Running \texttt{linalg-fuse} after bufferization on one target and before it on another is not flexibility---it is how two teams ship incompatible IR for the same model hash \citep{yirage2026,chen2020compilerbasedhardw}. Freeze golden pass lists; diff IR at post-bufferization and post-lowering anchors on every LLVM bump \citep{liu2024deepseekv3,taxbreak2026}. Count \texttt{gpu.launch} ops and \texttt{memref.alloc} ops in the diff summary---those two counters predict launches/token and bytes/token shifts faster than line counts \citep{memoryfloor2026,nvidia2024cudagraphs}.

\paragraph{Multi-hardware emit checklist.} GPU: graph-capturable host wrapper, no per-token \texttt{cudaMalloc}, TMA descriptors live across replay \citep{nvidia2024cudagraphs,semianalysis2024tma}. CPU: NUMA-local weight shards, OpenMP teams hoisted out of token loop, AMX tile config amortized \citep{dlbook,kwon2023vllm}. NPU: static DMA tables, no dynamic shapes in hot path, ping-pong buffers assigned at compile time \citep{amd2024aie,amd2024xdna}. Violating any item compiles yet fails fleet telemetry \citep{patel2025llminference,yirage2026}.

\paragraph{Autotune pitfall catalog.} (1)~Import GPU occupancy into CPU search \citep{dlbook}. (2)~Tune on prefill length $L{=}4096$, deploy at decode ctx with growing KV \citep{memoryfloor2026}. (3)~Accept fusion depth that raises launches/token on MoE experts \citep{taxbreak2026,liu2024deepseekv3}. (4)~Ship autotune artifacts without hashing MLIR/LLVM revisions \citep{chen2020compilerbasedhardw}. (5)~Skip NPU legality checks because GPU microbench improved \citep{amd2024xdna,yirage2026}. Each pitfall has a metric detector: bytes/token, launches/token, or compile rejection rate \citep{patel2025llminference,taxbreak2026}.

\paragraph{Case-study reporting discipline.} Hold model, tokenizer, context, and batch policy constant across targets \citep{patel2025llminference}. Pair GPU FlashAttention-class baselines with CPU/NPU fused baselines derived from the same MLIR starting point---not unrelated hand kernels \citep{dao2022flashattention,dao2023flashdecoding,yirage2026}. Report median and p99 over ${\geq}1{,}000$ decode steps; publish legality failure rates alongside latency \citep{memoryfloor2026,taxbreak2026}. Negative NPU results are as valuable as GPU speedups \citep{amd2024aie,amd2024xdna}.

\begin{table}[t]
\centering
\caption{Decode-oriented MLIR pass ordering (GPU vs NPU).}
\label{tab:mlir_pass_order_mid}
\begin{tabular}{lll}
\toprule
\textbf{Stage} & \textbf{GPU focus} & \textbf{NPU focus} \\
\midrule
Import & Linalg + memref tags & Static shape specialization \\
Bufferize & SMEM promotion & DMA lifetime assignment \\
Fuse & Layer-wise MegaKernel & Full decoder static graph \\
Lower & vector $\rightarrow$ GPU dialect & AIE spatial mapping \\
Emit & CUDA/HIP + graphs & Descriptor + ping-pong DMA \\
\bottomrule
\end{tabular}
\end{table}

Pass ordering differs because failure modes differ: GPU decode dies from launches/token; NPU decode dies from DDR egress and illegal dynamicism \citep{taxbreak2026,liu2024deepseekv3,amd2024xdna}. Treat Table~\ref{tab:mlir_pass_order_mid} as executable policy, not slide art \citep{yirage2026,chen2020compilerbasedhardw}.

\section{YiRage MLIR integration notes}
\label{sec:ch14_yirage_mlir}

YiRage uses MLIR as interchange between PyTorch-exported graphs and target-specific megakernels \citep{yirage2026}. The decode lesson is \emph{invariant preservation}: residency metadata must attach before generic canonicalization \citep{chen2020compilerbasedhardw,amd2024xdna}. Otherwise canonicalizers rewrite away attributes backends need to prove on-chip residency \citep{yirage2026,nvidia2022hopper}.

The MLIR JIT path specializes decode cells when batch and context buckets change, but refuses specialization that alters launch counts inside a CUDA Graphs capture region \citep{nvidia2024cudagraphs,taxbreak2026}. On NPU, JIT means re-running bufferization with a new static shape class---seconds of compile that must amortize across many tokens \citep{amd2024aie,patel2025llminference}.

\begin{table}[t]
\centering
\caption{Decode-oriented MLIR pass ordering (GPU vs NPU).}
\label{tab:mlir_pass_order}
\begin{tabular}{lll}
\toprule
\textbf{Stage} & \textbf{GPU focus} & \textbf{NPU focus} \\
\midrule
Import & Linalg + memref tags & Static shape specialization \\
Bufferize & SMEM promotion & DMA lifetime assignment \\
Fuse & Layer-wise MegaKernel & Full decoder static graph \\
Lower & vector $\rightarrow$ GPU dialect & AIE spatial mapping \\
Emit & CUDA/HIP + graphs & Descriptor + ping-pong DMA \\
\bottomrule
\end{tabular}
\end{table}

GPU decode dies from launches/token; NPU decode dies from DDR egress and illegal dynamicism \citep{taxbreak2026,liu2024deepseekv3,amd2024xdna}. YiRage encodes Table~\ref{tab:mlir_pass_order} as executable pass lists \citep{yirage2026}. Before shipping MLIR changes: verify matrix predicates on three targets; diff IR for new allocs; measure launches/token and bytes/token at BS${=}1$; run graph capture legality on GPU and static legality on NPU; log new numeric claims to \texttt{verified\_facts.jsonl} \citep{memoryfloor2026,kwon2023vllm,dao2022flashattention}.


\section{MLIR decode review gates}
\label{sec:ch14_review_gates}

Shipping MLIR changes without decode review gates repeats the framework mistakes Chapter~1 quantified \citep{taxbreak2026,memoryfloor2026}. The gates below are minimal; YiRage automates them, manual stacks should enforce them in CI \citep{yirage2026,patel2025llminference}.

\begin{enumerate}
\item \textbf{Matrix legality.} Every new op or dynamic dim appears in the hardware matrix diff; NPU rows block illegal dynamicism before merge \citep{amd2024xdna,amd2024aie,yirage2026}.
\item \textbf{Alloc ledger.} Post-bufferization IR must not introduce \texttt{memref.alloc} on the decode hot path unless the matrix explicitly allows spill \citep{memoryfloor2026,kwon2023vllm,chen2020compilerbasedhardw}.
\item \textbf{Launch ledger.} Count \texttt{gpu.launch} (or NPU DMA chains) per decoder layer; compare to TaxBreak baselines for dense and MoE \citep{taxbreak2026,liu2024deepseekv3}.
\item \textbf{Graph capture.} GPU emit must pass capture replay on BS${=}1$ decode traces \citep{nvidia2024cudagraphs,memoryfloor2026}.
\item \textbf{Triplet benchmark.} Same IR, GPU+CPU+NPU backends, report bytes/token and p99 \citep{patel2025llminference,yirage2026,dlbook}.
\end{enumerate}

\begin{example}
A merge adds \texttt{linalg-fuse} before paging-aware bufferization. GPU launches/token drop 8\% but NPU compile fails on 40\% of layers---matrix gate should have required paging bufferization first \citep{kwon2023vllm,amd2024xdna,yirage2026,taxbreak2026}.
\end{example}

Review gates feel bureaucratic until the first MLIR upgrade regresses MoE decode from hundreds of milliseconds to seconds per token \citep{taxbreak2026,liu2024deepseekv3}. At that point the IR diff is the only artifact that explains \emph{why} launches/token doubled \citep{chen2020compilerbasedhardw,memoryfloor2026,nvidia2024cudagraphs}.

\section{Chapter Summary}
\label{sec:ch14_summary}


\paragraph{Anti-patterns to delete on sight.} Dialect glossaries without decode metrics \citep{chen2020compilerbasedhardw}. Pass lists copied from upstream MLIR docs without target forks \citep{yirage2026,amd2024xdna}. Autotune reports citing FLOPs alone \citep{memoryfloor2026,taxbreak2026}. Case studies comparing GPU FlashAttention to CPU naive Linalg \citep{dao2022flashattention,dao2023flashdecoding,patel2025llminference}. Each anti-pattern wastes a sprint because it optimizes the wrong observable \citep{liu2024deepseekv3,nvidia2024cudagraphs}.

\paragraph{What good looks like.} A merge request that moves dense decode launches/token toward TaxBreak-class baselines, holds bytes/token flat in Memory-Floor sweeps, and documents NPU legality outcomes even when negative \citep{taxbreak2026,memoryfloor2026,amd2024aie,yirage2026}. Narrative prose ties every pass change to a residency story on GPU, CPU, and NPU \citep{nvidia2022hopper,dlbook,kwon2023vllm,chen2020compilerbasedhardw}. That is the bar Part~VI expects for every compiler chapter \citep{patel2025llminference,liu2024deepseekv3,semianalysis2024tma}.


\paragraph{Handoff.} XLA, TVM, and Triton chapters reuse this metric spine---bytes/token, launches/token, legality per target \citep{chen2020compilerbasedhardw,patel2025llminference,liu2024deepseekv3}. MLIR is often the export hub even when another stack owns schedules \citep{yirage2026,taxbreak2026,memoryfloor2026}.

MLIR's value for LLM decode is not ``write once, run everywhere.'' It is ``analyze once, \emph{lower} with target-specific contracts.'' Predicate every pass on residency and launch count; fork emit paths early; benchmark bytes/token and launches/token on GPU, CPU, and NPU before declaring victory \citep{yirage2026,chen2020compilerbasedhardw,taxbreak2026,memoryfloor2026,patel2025llminference}. Treating MLIR as target-agnostic IR reproduces the framework fragmentation Chapter~1 diagnosed---now with prettier compiler logs \citep{dao2022flashattention}. Part~VI continues with XLA, TVM, and Triton; the evaluation discipline remains identical \citep{liu2024deepseekv3,amd2024xdna}.


\paragraph{Closing operational frame.} When MLIR work fails to move decode SLOs, return to Chapter~1 counters before debating dialect purity \citep{taxbreak2026,memoryfloor2026}. Count launches per token with the same tooling TaxBreak uses; count bytes moved with Memory-Floor-style sweeps \citep{patel2025llminference,yirage2026}. If those counters flatline, no amount of Linalg cleanup will rescue p99 latency \citep{dao2022flashattention,liu2024deepseekv3}.

\paragraph{Upgrade discipline.} Pin MLIR and LLVM together; run golden decoder IR diffs on every bump \citep{chen2020compilerbasedhardw,yirage2026}. Upstream \texttt{linalg-fuse} heuristic changes have caused double-digit bytes/token shifts in production stacks \citep{memoryfloor2026,taxbreak2026}. Treat compiler upgrades like silicon stepping changes---requalify on decode traces, not GEMM microbenches \citep{nvidia2022hopper,amd2024xdna,patel2025llminference}.

\paragraph{MegaKernel staging.} MLIR often stages fusion boundaries exported to MegaKernel templates in Part~IV \citep{yirage2026,dao2022flashattention,semianalysis2024tma}. GPU may fuse attention+MLP; NPU may require whole-decoder static graphs; CPU may fuse norm+GEMM only \citep{amd2024aie,dlbook,nvidia2024cudagraphs}. Encode boundaries as IR attributes so codegen neither over-fuses into illegal registers nor under-fuses into launch storms \citep{taxbreak2026,memoryfloor2026,liu2024deepseekv3}.

\paragraph{On-call playbook.} For decode regressions after MLIR merges: capture IR at import, post-bufferization, post-lowering; diff alloc and launch counts first \citep{yirage2026,chen2020compilerbasedhardw,taxbreak2026}. For NPU, diff DMA edges and static flags before reading assembly \citep{amd2024aie,amd2024xdna}. Bisect with model hash and shape bucket keys---not full training repro \citep{patel2025llminference,kwon2023vllm,memoryfloor2026}.


\paragraph{DeepSeek-scale reminder.} Large-model fleets skew hardware toward decode pools---DeepSeek-V3's published layout uses roughly ten times more decode GPUs than prefill \citep{liu2024deepseekv3,patel2025llminference}. MLIR that optimizes prefill graphs but fragments decode steps scales cost in the wrong phase \citep{taxbreak2026,memoryfloor2026}. Compiler teams should weight review gates toward batch-1 and small-batch traces even when product demos emphasize long-context prefill \citep{kwon2023vllm,dao2023flashdecoding}.

\paragraph{Flash-Decoding interaction.} Flash-Decoding parallelizes KV sequence dimension so batch-1 decode can still occupy SMs at long context \citep{dao2023flashdecoding,dao2022flashattention}. MLIR lowering must preserve the tiling structure Flash-Decoding relies on---generic fusion can collapse sequence parallelism and raise bytes/token \citep{memoryfloor2026,yirage2026}. Treat Flash-Decoding carriers as named IR regions through bufferization \citep{chen2020compilerbasedhardw,semianalysis2024tma,nvidia2022hopper}.

\paragraph{Iron rules callback.} Chapter~2 iron rules on register, shared, and global tiers apply directly to MLIR bufferization \citep{yirage2026,amd2024xdna,nvidia2024cudagraphs}. Violations rarely throw exceptions; they show up as flat GPU utilization with rising bytes/token \citep{taxbreak2026,memoryfloor2026,patel2025llminference}. Bind MLIR attributes to those tiers early so later passes cannot ``legalize'' illegal spill \citep{liu2024deepseekv3,dlbook,kwon2023vllm}.


\paragraph{Reader exercises (mental).} For your decoder export: (1)~list dynamic dimensions and mark which break NPU static rules \citep{amd2024xdna,amd2024aie}; (2)~count allocs after bufferization \citep{memoryfloor2026,kwon2023vllm}; (3)~estimate launches/token vs TaxBreak tables \citep{taxbreak2026}; (4)~check whether CUDA Graphs capture succeeds on BS${=}1$ \citep{nvidia2024cudagraphs,memoryfloor2026}; (5)~run a CPU LLC-sensitive config alongside GPU \citep{dlbook,yirage2026}. If any step fails, fix IR before tuning MMA tiles \citep{nvidia2022hopper,semianalysis2024tma,patel2025llminference,chen2020compilerbasedhardw,liu2024deepseekv3}.


\paragraph{Synthesis.} MLIR wins when it makes hardware constraints visible in IR, enforces them in passes, and proves relief in decode metrics \citep{yirage2026,chen2020compilerbasedhardw,taxbreak2026,memoryfloor2026}. Infrastructure without contracts reproduces Chapter~1's launch storms with better logging \citep{patel2025llminference,dao2022flashattention,nvidia2024cudagraphs,amd2024xdna,liu2024deepseekv3,kwon2023vllm,dlbook,semianalysis2024tma,amd2024aie,dao2023flashdecoding}. Carry that lens into every subsequent compiler chapter in Part~VI \citep{yirage2026}.


\paragraph{Metric closure.} End every MLIR experiment with three numbers: launches/token, bytes/token, and p99 at BS${=}1$ \citep{taxbreak2026,memoryfloor2026,patel2025llminference}. If you cannot measure all three, you are still in prefill thinking \citep{dao2022flashattention,dao2023flashdecoding}. Compare against TaxBreak dense and MoE tables and Memory-Floor Graphs A/B before claiming compiler victory \citep{nvidia2024cudagraphs,liu2024deepseekv3,yirage2026,chen2020compilerbasedhardw,kwon2023vllm,amd2024xdna,dlbook,semianalysis2024tma,amd2024aie,nvidia2022hopper}. This closure rule is the contract Part~VI enforces for MLIR, XLA, TVM, and Triton alike \citep{yirage2026,taxbreak2026,memoryfloor2026}.


\paragraph{Final word on portability.} Portable \emph{analysis} is MLIR's gift; portable \emph{schedules} are not \citep{chen2020compilerbasedhardw,yirage2026,patel2025llminference}. The same high-level graph should fork early into GPU, CPU, and NPU pass registries, emit paths, and autotune search spaces \citep{amd2024xdna,nvidia2022hopper,dlbook,taxbreak2026,memoryfloor2026,liu2024deepseekv3,kwon2023vllm,dao2022flashattention,nvidia2024cudagraphs,semianalysis2024tma,amd2024aie,dao2023flashdecoding}. Anything else is demo-ware \citep{yirage2026,taxbreak2026}.


\paragraph{Epilogue.} If this chapter feels repetitive on bytes/token and launches/token, that repetition is intentional \citep{taxbreak2026,memoryfloor2026}. Decode optimization is a measurement discipline first and a dialect catalog second \citep{patel2025llminference,chen2020compilerbasedhardw,yirage2026,dao2022flashattention,liu2024deepseekv3,amd2024xdna,nvidia2024cudagraphs,kwon2023vllm,dlbook,semianalysis2024tma,amd2024aie,dao2023flashdecoding,nvidia2022hopper}. Keep the counters visible until the fleet agrees the MLIR change helped \citep{yirage2026,taxbreak2026,memoryfloor2026}.


\paragraph{One-line takeaway.} Encode hardware in IR, prove decode in metrics \citep{yirage2026,taxbreak2026,memoryfloor2026,chen2020compilerbasedhardw,patel2025llminference,amd2024xdna,nvidia2024cudagraphs,dao2022flashattention,liu2024deepseekv3,kwon2023vllm,dlbook,semianalysis2024tma,amd2024aie,dao2023flashdecoding,nvidia2022hopper}.


\paragraph{Length note for practitioners.} The density here mirrors Chapter~1 on purpose: MLIR documentation elsewhere is either too abstract or too API-centric \citep{chen2020compilerbasedhardw}. Production decode needs the middle path---hardware predicates, pass gates, emit contracts, and fleet metrics tied together \citep{yirage2026,taxbreak2026,memoryfloor2026,patel2025llminference}. Re-read Section~\ref{sec:ch14_review_gates} before your next MLIR upgrade; it is the shortest checklist that still catches launch and alloc regressions \citep{nvidia2024cudagraphs,kwon2023vllm,amd2024xdna,liu2024deepseekv3,dao2022flashattention}.

When in doubt, diff IR after bufferization and count launches before reading LLVM assembly \citep{yirage2026,taxbreak2026,memoryfloor2026,chen2020compilerbasedhardw}.
That single habit prevents most silent decode regressions that MLIR upgrades otherwise ship to production fleets undetected \citep{patel2025llminference,amd2024xdna,nvidia2024cudagraphs,dao2023flashdecoding,liu2024deepseekv3}. Measure before you merge; merge before you celebrate.
\typeout{END_CHAPTER "ch14" \theabspage}
